\subsection{La piattaforma Serverledge}
\label{subsec:serverledge}

Serverledge è una piattaforma Function-as-a-Service progettata
per ambienti cloud-edge, caratterizzata da un'architettura
decentralizzata che consente l'esecuzione di funzioni serverless
all'interno di container. La piattaforma separa le componenti di
controllo da quelle di esecuzione, permettendo una gestione
distribuita delle istanze e delle informazioni di stato tra i
nodi che compongono l'infrastruttura.

Ogni nodo Serverledge espone un'interfaccia per la creazione,
l'eliminazione e l'invocazione delle funzioni, supportando sia
la modalità sincrona sia quella asincrona. Le decisioni di
esecuzione sono affidate a uno scheduler, responsabile del
\emph{placement} delle richieste sui nodi disponibili.
L'esecuzione avviene tramite container Docker, mantenuti attivi
per un intervallo di tempo limitato successivamente all'invocazione,
con l'obiettivo di ridurre l'impatto dei \emph{cold start} sulle
invocazioni successive.

La condivisione delle informazioni tra nodi è realizzata tramite
un meccanismo di storage distribuito, che consente di propagare
metadati e informazioni di stato all'interno del cluster. Tali
informazioni sono accessibili dallo scheduler al momento della
decisione di \emph{placement}, permettendo di scegliere tra
l'esecuzione locale sul nodo ricevente e l'offloading verso
altri nodi edge o verso il cloud, in funzione delle condizioni
operative rilevate.

Dal punto di vista architetturale, Serverledge è progettata per
essere estensibile: le logiche di scheduling e di gestione delle
funzioni sono separate dal modello di interazione lato utente,
rendendo possibile l'integrazione di componenti aggiuntivi di
monitoraggio e controllo senza modificare l'interfaccia esposta
agli utenti. Questa proprietà è particolarmente rilevante nel
contesto del presente lavoro, in quanto consente di introdurre
meccanismi di gestione energetica direttamente a livello di
scheduler, sfruttando i punti di estensione nativi della
piattaforma senza alterarne il comportamento di base.